\documentclass[10pt,a4paper]{amsart}
\usepackage{amsmath,amssymb,amsthm,graphicx}
\usepackage{verbatim}
\newenvironment{reference}{\comment}{\endcomment}
%\newenvironment{reference}{}{}
\newenvironment{slogan}{\comment}{\endcomment}
\newenvironment{history}{\comment}{\endcomment}

% For commutative diagrams we use Xy-pic
\usepackage[all]{xy}

% We use 2cell for 2-commutative diagrams.
\xyoption{2cell}
% Theorem environments.
%
\theoremstyle{plain}
\newtheorem{theorem}[subsection]{Theorem}
\newtheorem{proposition}[subsection]{Proposition}
\newtheorem{lemma}[subsection]{Lemma}

\theoremstyle{definition}
\newtheorem{definition}[subsection]{Definition}
\newtheorem{example}[subsection]{Example}
\newtheorem{exercise}[subsection]{Exercise}
\newtheorem{situation}[subsection]{Situation}

\theoremstyle{remark}
\newtheorem{remark}[subsection]{Remark}
\newtheorem{remarks}[subsection]{Remarks}

\numberwithin{equation}{subsection}

% Macros
%
\def\lim{\mathop{\mathrm{lim}}\nolimits}
\def\colim{\mathop{\mathrm{colim}}\nolimits}
\def\Spec{\mathop{\mathrm{Spec}}}
\def\Hom{\mathop{\mathrm{Hom}}\nolimits}
\def\Ext{\mathop{\mathrm{Ext}}\nolimits}
\def\SheafHom{\mathop{\mathcal{H}\!\mathit{om}}\nolimits}
\def\SheafExt{\mathop{\mathcal{E}\!\mathit{xt}}\nolimits}
\def\Sch{\mathit{Sch}}
\def\Mor{\mathop{\mathrm{Mor}}\nolimits}
\def\Ob{\mathop{\mathrm{Ob}}\nolimits}
\def\Sh{\mathop{\mathit{Sh}}\nolimits}
\def\NL{\mathop{N\!L}\nolimits}
\def\CH{\mathop{\mathrm{CH}}\nolimits}
\def\proetale{{pro\text{-}\acute{e}tale}}
\def\etale{{\acute{e}tale}}
\def\QCoh{\mathit{QCoh}}
\def\Ker{\mathop{\mathrm{Ker}}}
\def\Im{\mathop{\mathrm{Im}}}
\def\Coker{\mathop{\mathrm{Coker}}}
\def\Coim{\mathop{\mathrm{Coim}}}

% Boxtimes
%
\DeclareMathSymbol{\boxtimes}{\mathbin}{AMSa}{"02}

%
% Macros for moduli stacks/spaces
%
\def\QCohstack{\mathcal{QC}\!\mathit{oh}}
\def\Cohstack{\mathcal{C}\!\mathit{oh}}
\def\Spacesstack{\mathcal{S}\!\mathit{paces}}
\def\Quotfunctor{\mathrm{Quot}}
\def\Hilbfunctor{\mathrm{Hilb}}
\def\Curvesstack{\mathcal{C}\!\mathit{urves}}
\def\Polarizedstack{\mathcal{P}\!\mathit{olarized}}
\def\Complexesstack{\mathcal{C}\!\mathit{omplexes}}
% \Pic is the operator that assigns to X its picard group, usage \Pic(X)
% \Picardstack_{X/B} denotes the Picard stack of X over B
% \Picardfunctor_{X/B} denotes the Picard functor of X over B
\def\Pic{\mathop{\mathrm{Pic}}\nolimits}
\def\Picardstack{\mathcal{P}\!\mathit{ic}}
\def\Picardfunctor{\mathrm{Pic}}
\def\Deformationcategory{\mathcal{D}\!\mathit{ef}}
\usepackage{cleveref}
\crefname{originallemma}{original lemma}{original lemmas}
\Crefname{originallemma}{Original lemma}{Original lemmas}
\makeatletter
\newlabel{lemma-localize-morphism}{{03CF}{source}{}{sourcepointer-1}{}}
\newlabel{lemma-localize-morphism@cref}{{[originallemma][0][]03CF}{[1][0][]source}{}{}{}}
\makeatother
\crefname{originallemma}{original lemma}{original lemmas}
\Crefname{originallemma}{Original lemma}{Original lemmas}
\makeatletter
\newlabel{lemma-relocalize}{{03EH}{source}{}{sourcepointer-2}{}}
\newlabel{lemma-relocalize@cref}{{[originallemma][0][]03EH}{[1][0][]source}{}{}{}}
\makeatother
\crefname{originallemma}{original lemma}{original lemmas}
\Crefname{originallemma}{Original lemma}{Original lemmas}
\makeatletter
\newlabel{lemma-essential-image-j-shriek}{{00Y1}{source}{}{sourcepointer-3}{}}
\newlabel{lemma-essential-image-j-shriek@cref}{{[originallemma][0][]00Y1}{[1][0][]source}{}{}{}}
\makeatother
\crefname{originalequation}{original equation}{original equations}
\Crefname{originalequation}{Original equation}{Original equations}
\makeatletter
\newlabel{equation-relocalize}{{04IK}{source}{}{sourcepointer-4}{}}
\newlabel{equation-relocalize@cref}{{[originalequation][0][]04IK}{[1][0][]source}{}{}{}}
\makeatother
\crefname{originallemma}{original lemma}{original lemmas}
\Crefname{originallemma}{Original lemma}{Original lemmas}
\makeatletter
\newlabel{lemma-relocalize-explicit}{{04IL}{source}{}{sourcepointer-5}{}}
\newlabel{lemma-relocalize-explicit@cref}{{[originallemma][0][]04IL}{[1][0][]source}{}{}{}}
\makeatother
\crefname{originallemma}{original lemma}{original lemmas}
\Crefname{originallemma}{Original lemma}{Original lemmas}
\makeatletter
\newlabel{lemma-pullback-representable-sheaf}{{04D3}{source}{}{sourcepointer-6}{}}
\newlabel{lemma-pullback-representable-sheaf@cref}{{[originallemma][0][]04D3}{[1][0][]source}{}{}{}}
\makeatother
\title{The localization of a morphism of sites along a morphism of objects}
\author{The Stacks Project Authors}
\date{stacks-project:sites.tex:local-snapshot-20261012; GNU FDL 1.2 or later}
\begin{document}
\maketitle
\noindent\textit{Editorial note.} All delivered mathematical text is the author's own, copied from the pinned Stacks Project chapter file sites-cohomology.tex; no statement or proof step was rewritten, reordered or removed.\par\medskip
\noindent\textit{Editorial note.} The proof is the author's argument and is explicitly a collection of the identifications available at this point of the development: each of the three assertions is reduced to a named result or to a displayed formal identity, and the author's words describing it as such are preserved verbatim inside the counted span.\par\medskip
\noindent\textit{Editorial note.} Every labelled reference occurring inside the counted spans is carried in assembly.external\_references as a non-load-bearing navigation pointer (same-chapter ones at their definition site, cross-chapter ones at the invocation inside the counted span), so no load-bearing mathematics is replaced by a stub. No bibliography citation occurs inside the counted spans.\par\medskip
\noindent\textit{Editorial note.} The checked chapter list for same-name collisions is the set of stacks chapters staged in this repository: chow, cohomology, divisors, etale-cohomology, flat, more-algebra, more-morphisms, morphisms, spaces-more-morphisms, varieties together with this card's chapter sites-cohomology.\par\medskip
\noindent\textit{Editorial note.} The preamble is the chapter's own retained minimal preamble (sources/stacks-sites-cohomology/retained-preamble.tex), a verbatim excerpt of the native preamble: the Xy-pic setup and the native theorem declarations with the mathematical macros.\par\medskip
\noindent\textit{Allowed standard prerequisites (editorial).}\begin{itemize}
\item The localization of a morphism of sites. Same chapter, Lemma lemma-localize-morphism. Supplies the morphism between the localized sites and its description.
\item Relocalizing a site. Same chapter, Lemma lemma-relocalize and Lemma lemma-relocalize-explicit. Supply the composite factorization of the morphism and the explicit description of the pullback.
\item The pullback of a representable sheaf. Same chapter, Lemma lemma-pullback-representable-sheaf. Supplies the identifications used in the formal identity of fibre products.
\end{itemize}\medskip
\ifdefined\hypertarget\hypertarget{sourcepointer-1}{}\fi
\noindent\textit{Source reference.} Original lemma 03CF: sources/stacks-sites/pinned-original-source.tex, source lines 6601--6601. Navigation pointer invoked at proof line 6601; the counted span contains the pointer only and the cited result is not reproduced inside it, so the pointer is registered as a non-load-bearing navigation stub in the established form for Stacks material. No bibliography citation is involved.\par\medskip
\ifdefined\hypertarget\hypertarget{sourcepointer-2}{}\fi
\noindent\textit{Source reference.} Original lemma 03EH: sources/stacks-sites/pinned-original-source.tex, source lines 6681--6681. Navigation pointer invoked at proof line 6681; the counted span contains the pointer only and the cited result is not reproduced inside it, so the pointer is registered as a non-load-bearing navigation stub in the established form for Stacks material. No bibliography citation is involved.\par\medskip
\ifdefined\hypertarget\hypertarget{sourcepointer-3}{}\fi
\noindent\textit{Source reference.} Original lemma 00Y1: sources/stacks-sites/pinned-original-source.tex, source lines 6685--6685. Navigation pointer invoked at proof line 6685; the counted span contains the pointer only and the cited result is not reproduced inside it, so the pointer is registered as a non-load-bearing navigation stub in the established form for Stacks material. No bibliography citation is involved.\par\medskip
\ifdefined\hypertarget\hypertarget{sourcepointer-4}{}\fi
\noindent\textit{Source reference.} Original equation 04IK: sources/stacks-sites/pinned-original-source.tex, source lines 6717--6717. Navigation pointer invoked at proof line 6717; the counted span contains the pointer only and the cited result is not reproduced inside it, so the pointer is registered as a non-load-bearing navigation stub in the established form for Stacks material. No bibliography citation is involved.\par\medskip
\ifdefined\hypertarget\hypertarget{sourcepointer-5}{}\fi
\noindent\textit{Source reference.} Original lemma 04IL: sources/stacks-sites/pinned-original-source.tex, source lines 6721--6721. Navigation pointer invoked at proof line 6721; the counted span contains the pointer only and the cited result is not reproduced inside it, so the pointer is registered as a non-load-bearing navigation stub in the established form for Stacks material. No bibliography citation is involved.\par\medskip
\ifdefined\hypertarget\hypertarget{sourcepointer-6}{}\fi
\noindent\textit{Source reference.} Original lemma 04D3: sources/stacks-sites/pinned-original-source.tex, source lines 6734--6734. Navigation pointer invoked at proof line 6734; the counted span contains the pointer only and the cited result is not reproduced inside it, so the pointer is registered as a non-load-bearing navigation stub in the established form for Stacks material. No bibliography citation is involved.\par\medskip
\begin{lemma}
\label{lemma-relocalize-morphism}
Let $f : \mathcal{C} \to \mathcal{D}$ be a morphism of sites
corresponding to the continuous functor $u : \mathcal{D} \to \mathcal{C}$.
Let $V \in \Ob(\mathcal{D})$, $U \in \Ob(\mathcal{C})$
and $c : U \to u(V)$ a morphism of $\mathcal{C}$.
There exists a commutative diagram of topoi
$$
\xymatrix{
\Sh(\mathcal{C}/U) \ar[r]_{j_U} \ar[d]_{f_c} &
\Sh(\mathcal{C}) \ar[d]^f \\
\Sh(\mathcal{D}/V) \ar[r]^{j_V} &
\Sh(\mathcal{D}).
}
$$
We have $f_c = f' \circ j_{U/u(V)}$ where
$f' : \Sh(\mathcal{C}/u(V)) \to \Sh(\mathcal{D}/V)$
is as in
Lemma \ref{lemma-localize-morphism}
and
$j_{U/u(V)} : \Sh(\mathcal{C}/U) \to \Sh(\mathcal{C}/u(V))$
is as in
Lemma \ref{lemma-relocalize}.
Using the identifications
$\Sh(\mathcal{C}/U) = \Sh(\mathcal{C})/h_U^\#$ and
$\Sh(\mathcal{D}/V) = \Sh(\mathcal{D})/h_V^\#$ of
Lemma \ref{lemma-essential-image-j-shriek}
the functor $(f_c)^{-1}$ is described by the rule
$$
(f_c)^{-1}(\mathcal{H} \xrightarrow{\varphi} h_V^\#)
=
(f^{-1}\mathcal{H} \times_{f^{-1}\varphi, h_{u(V)}^\#, c} h_U^\#
\rightarrow h_U^\#).
$$
Finally, given any morphisms $b : V' \to V$, $a : U' \to U$ and
$c' : U' \to u(V')$ such that
$$
\xymatrix{
U' \ar[r]_-{c'} \ar[d]_a & u(V') \ar[d]^{u(b)} \\
U \ar[r]^-c & u(V)
}
$$
commutes, then the diagram
$$
\xymatrix{
\Sh(\mathcal{C}/U') \ar[r]_{j_{U'/U}} \ar[d]_{f_{c'}} &
\Sh(\mathcal{C}/U) \ar[d]^{f_c} \\
\Sh(\mathcal{D}/V') \ar[r]^{j_{V'/V}} &
\Sh(\mathcal{D}/V).
}
$$
commutes.
\end{lemma}

\begin{proof}
This lemma proves itself, and is more a collection of things we know
at this stage of the development of theory. For example the commutativity
of the first square follows from the commutativity of
Diagram (\ref{equation-relocalize})
and the commutativity of the diagram in
Lemma \ref{lemma-localize-morphism}.
The description of $f_c^{-1}$ follows on combining
Lemma \ref{lemma-relocalize-explicit}
with
Lemma \ref{lemma-localize-morphism}.
The commutativity of the last square then follows from the
equality
$$
f^{-1}\mathcal{H} \times_{h_{u(V)}^\#, c} h_U^\# \times_{h_U^\#} h_{U'}^\#
=
f^{-1}(\mathcal{H} \times_{h_V^\#} h_{V'}^\#)
\times_{h_{u(V'), c'}^\#} h_{U'}^\#
$$
which is formal using that $f^{-1}h_V^\# = h_{u(V)}^\#$ and
$f^{-1}h_{V'}^\# = h_{u(V')}^\#$, see
Lemma \ref{lemma-pullback-representable-sheaf}.
\end{proof}

\end{document}
